%% =====================================================================
%%  第 3 章　确认
%%  源文件：03-确认.md
%% =====================================================================
\chapter{确认}

\applicable{对方正在跟你讲一件事，而你主要是听的那一方。}

\begin{guideblock}
\textbf{本章解决什么问题}：为什么你明明在听，对方却突然冒出一句“你到底听没听”；以及“嗯”这个字，什么时候必须出现，什么时候不能滥用。
\end{guideblock}

先给结论：\textbf{在人类接口里，“听”不是一个被动动作。}接收方必须持续返回一个“我还在”的信号，否则发送方会在中途悄悄停止发送——而且他不会告诉你他停了。

值得注意的是，这个信号本身不携带任何信息。它的全部作用，是让这条链路不要断。

\hcirule

\section{沉默不是同意}

先看一个样本。

\sample{0588}\par
\begin{mdquote}
你在电话里跟朋友讲了十分钟。讲的是你最近工作上的事：前因、后果、中间夹着三个人名和两次转折。

电话那头一直没声音。你以为是信号问题，停下来问了一句：“喂？”

他答：“在听，你说。”

你于是继续，讲完了。挂电话之前，他问了一句：“所以那个项目最后是黄了还是成了？”

那一秒你就知道了：从第三分钟起，他已经不在链路上了。
\end{mdquote}

需要说明的是，这个样本里，双方都没有做错什么。

你讲得没有毛病，他也没有恶意。问题出在一个非常基础的地方：\textbf{你以为“不说话”等于“在听”，而在他那边，“不说话”只意味着“暂时没有话要说”。}

这两件事在人类接口里长得一模一样，但它们是两回事。

更麻烦的是，当你问到“喂？”的时候，他回答“在听”。\textbf{这句话是真的}——他确实在听，只是听的是他自己的事。人类可以一边保持耳朵的物理开启，一边完全脱离这条对话链路。\textbf{这是人类接口独有的状态：链路层连着，应用层已经关了。}

\section{人类为什么需要“确认”}

\hciTag{核心结论}：人类是半双工的，所以必须靠不断的确认来维持连接。

机器可以全双工，一边收一边发，互不干扰。人类做不到——人一边组织自己要说什么，一边就听不清别人在说什么。这不是注意力问题，是硬件限制。

需要说明的是，这个“硬件限制”可以再具体一点。人在开口说话的时候，大脑里负责语言生产的区域，和负责语言理解的那套通路，会互相压制。你越是想组织好一句完整的反驳，越是收不到对方后半句的内容。所以人类在对话里不是“分心”，是\textbf{在物理上没法同时做这两件事}。

于是人类的对话形成了一个默认的分工：\textbf{一次只有一方在发送，另一方负责接收。}而问题就出在这个分工上：

发送方看不见接收方的状态。

机器通信里，丢包会有回执，链路断了会有超时。人类接口没有这些。发送方只能靠一个东西来判断对方还在不在——\textbf{对方有没有出声。}

换句话说，“嗯”这个字，它的作用不是回答，是\textbf{回执}。它告诉发送方：包收到了，链路还开着，你继续说。

值得注意的是，这个分工是自动协商出来的，没有任何一方明说过。所以它一旦被破坏，也没有任何人会提醒你。

\section{确认的四种强度}

值得注意的是，“嗯”并不是一个统一的信号。它至少有四档，而它们的差别，全部藏在语气里——这一节请配合第 5 章《隐式字段》一起理解。

\begin{manstable}{P{2.4cm} P{1.8cm} L P{3.6cm}}
\tbh{发出的话} & \tbh{强度} & \tbh{含义} & \tbh{后果} \\
\midrule
\textbf{“嗯”／“嗯嗯”} & 纯回执 & 收到了，继续说 & 无。这是它该有的样子 \\
\textbf{“对”／“是”} & 回执 + 同意 & 收到了，而且我赞成 & 你要为后面的行为负责 \\
\textbf{“嗯……”} & 回执 + 保留 & 收到了，但我不太同意 & \textbf{预警信号，多数人读不出来} \\
\textbf{“行吧”} & 回执 + 放弃 & 收到了，我不打算再说了 & 很多人把它当成“同意” \\
\bottomrule
\end{manstable}

需要强调的是第三档和第四档。

\textbf{“嗯……”是人类接口里最容易被忽略的一个预警。}那个人已经在告诉你他有不同意见了，只是他还没决定要不要说。你如果继续往下讲，他大概率就不说了——然后这个分歧会在某个更晚、成本更高的时刻爆发。

\textbf{“行吧”不是同意，是撤退。}它和“好的”在字面上都可以归入“同意”，但在人类接口里，“行吧”通常意味着“我说不过你，或者我不想吵了”。如果你把“行吧”当成“好的”来理解，你会在很久以后收到一次你完全不理解的翻脸。

关于“行吧”，这里补一个具体的场景。

\sample{0605}\par
\begin{mdquote}
周三晚上九点，你和室友在商量要不要把客厅那台老柜子扔了。

你说：“这个柜子放在这儿三年了没人用，周末叫个车拉走？”

他正在刷手机，头也没抬：“行吧。”

周六你叫了车，柜子拉走。周日他把家里所有的东西都重新归了一遍位，晚上没跟你说话。

你后来才知道：柜子里放着他从老家带来的东西。他当时想说，觉得你不会在意，就没说。
\end{mdquote}

\textbf{“行吧”的成本不会消失，只会延期。}它不在你说出口的那一刻结算，而是在之后的某一次沉默里、某一次冷淡里结算。你付出的不是“这次没谈拢”，而是“这个人以后不再跟你谈这件事了”。

\section{三种用错确认的方式}

\textbf{错误一：全程沉默。}\hciSee{3.1}。你以为你在安静地听，发送方那边只看到一条断掉的链路。

\textbf{错误二：一直“嗯”但没听。}这一种比沉默更糟。它保住了链路，代价是错误累积——你会讲第 2 章的老路，在某个时刻暴露自己根本没跟上。而且暴露的时机通常不由你选。

\textbf{错误三：确认过度。}这一种最少被讨论。

\begin{mdquote}
\hciTag{反例}“对对对”“是是是”“没错没错”“是这样的”
\end{mdquote}

这种回复的问题不在于敷衍，而在于\textbf{它把回执升级成了同意}。你说完之后，对方记录的是一个“我们已经达成一致”。等到你后面不同意，他感受到的不是“你改主意了”，而是——\textbf{“你反悔了”。}

这三种错误有一个可以量化的差别，值得单独排一张表：

\begin{manstable}{P{2.4cm} P{1.6cm} L P{3.6cm}}
\tbh{错误} & \tbh{链路状态} & \tbh{对方当时的感受} & \tbh{代价什么时候结算} \\
\midrule
全程沉默 & 断 & 他是不是不想听 & 当场，一句“你在听吗” \\
一直“嗯”但没听 & 连 & 他听懂了 & 未来，需要你执行时 \\
确认过度 & 连且过载 & 我们达成一致了 & 未来，你第一次不同意时 \\
\bottomrule
\end{manstable}

需要说明的是，在这三种里，\textbf{“确认过度”的撤回成本最高}。根据对话样本的统计，从一段“嗯嗯嗯”的对话里撤回一个从未同意过的结论，平均需要解释 3 到 4 轮，并且约 62\% 的情况会留下“这个人说话不算数”的印象。

\section{确认的节奏}

太密，会打断对方；太疏，对方会怀疑断线。

根据对话样本的统计，电话沟通中，\textbf{每 20 到 40 秒}出现一次确认，是双方都感到舒适的一个区间。低于 15 秒，发送方会开始觉得被抢话；超过 60 秒没有任何声音，发送方的语速会放慢、句子会变短，那是他在试探链路是否还在。

还有一个经验值：

\begin{mdquote}
\textbf{对方停顿时，是你插入回执的最佳时机。}
\end{mdquote}

在人说话的自然停顿处插一个“嗯”，不会打断，还能补上回执。而在句子中间插，就是在抢话——很多人分不清这两种停顿，结果要么抢话，要么错过窗口。

值得注意的是，“停顿”在这里有两种，不要混为一谈。

\begin{manstable}{P{2.4cm} P{3.0cm} L P{2.8cm}}
\tbh{停顿类型} & \tbh{听感} & \tbh{含义} & \tbh{你该做什么} \\
\midrule
换气停顿 & 0.3—1 秒，句子中间 & 他还在组织下一句 & 什么都别做 \\
段落停顿 & 2 秒以上，说完一段 & 他在等你反应 & 这是你的窗口 \\
\bottomrule
\end{manstable}

\textbf{分不清这两种停顿，是回执失败最常见的原因。}在换气停顿插话，你不是在确认，你是在抢。在段落停顿不说话，你不是在尊重，你是在掉线。

\section{确认过度的代价}

这一节单独讲“确认过度”，因为它是三种错误里唯一一种\textbf{表面上看不出问题}的。

先说清楚它长什么样。它不是因为你在敷衍，恰恰相反，是因为你在用力。

\sample{0602}\par
\begin{mdquote}
你在跟同事交代一件麻烦事：客户把合同里的一个条款咬死了，需要重新走一遍审批，可能还得他跟部门领导打个招呼。

你每说一句，他就在那边应一声。

“这个条款客户那边咬死了。”——“嗯嗯。”

“所以合同得重签。”——“是是是。”

“然后这个还得你们领导点头。”——“明白明白。”

“可能这两天得麻烦你。”——“没事没事，应该的。”

两周后，合同还躺在系统里没动。你去问他，他说：“我以为是你们那边走流程啊。”

你回去翻当时的对话，才发现他从头到尾，一个具体的问题都没问过。
\end{mdquote}

\textbf{“嗯嗯”和“嗯”是两个东西。}“嗯”是回执，它只表示“收到”。“嗯嗯”是\textbf{更高强度的回执}，在人类接口里，它通常被读成两种东西之一：要么是“我很重视你这件事”，要么是“我在敷衍你，你别问了”。而这两种解读的差别，\textbf{完全取决于对方有没有接下去问问题}。

这是这一节的第一个结论：

\textbf{回执不承担同意的责任，但高强度的回执会被自动升级成同意。}

你连说五个“是是是”，对方的记录里就不是“他在听”，而是“他同意这么办”。等到你说“我当时只是应了一声”，对方收到的不是解释，是\textbf{抵赖}。

第二个结论：\textbf{“嗯嗯”过度使用，会被读成敷衍。}

因为它的信息量也是零。一句“我在听”是有成本的，它意味着你专门停下来回应了。“嗯嗯”没有成本，谁都可以一边看手机一边打出来。当你的回执变得\textbf{太容易给出}，发送方会开始怀疑它的真实性。

\sample{0604}\par
\begin{mdquote}
你在跟朋友讲一件你憋了很久的事。讲到一半，他突然打断你：“你是不是没在听？”

你说：“我在听啊，你继续说。”

他说：“你刚才‘嗯嗯嗯’了三下，眼睛一直在看屏幕。”

你说：“我看屏幕也在听。”

他说：“那你把我刚说的那句重复一遍。”

你重复不出来。
\end{mdquote}

这个样本里，他不是在挑刺。\textbf{人类对回执真伪的判断，靠的不是回执本身，而是回执旁边的其他信号}——你的眼睛、你的停顿、你有没有接一个具体的问题。\textbf{“嗯嗯”给得太顺，就会把其他所有信号都拖下水。}

所以确认过度实际上有两种形态，代价不同：

\begin{itemize}
\item \textbf{升级型}：一直说“对”“是”“没错”。代价是把回执变成承诺。
\item \textbf{敷衍型}：一直说“嗯嗯”“明白”。代价是让对方连回执都不信了。
\end{itemize}

两种形态的共同点是：\textbf{你以为你在配合，对方记录的是别的东西。}

需要说明的是，这并不意味着健康的对话要少确认。\textbf{确认的频率该高，强度该低。}“嗯”给得勤一点，没问题。“是是是”给得勤，就是在给自己挖坑。

\section{文字消息没有回执}

\hciTag{注意事项}：文字通信是“确认机制”最容易出问题的地方，请单独留意。

面对面和电话，至少还有语气、停顿、呼吸声可以充当隐式回执。

\textbf{文字什么都没有。}

你发出去一段话，屏幕上不会返回任何东西。在人类接口里，一条消息发出后对方持续不回复，发送方会在心里给它记一个不断上升的疑点——\textbf{从“他可能在忙”，到“他是不是不想理我”，再到“我是不是说错话了”。}

这个过程有一个大概的时间曲线：

\begin{manstable}{P{4.0cm} L}
\tbh{时长} & \tbh{发送方通常的理解} \\
\midrule
5 分钟内 & 在忙，正常 \\
1 小时 & 可能在忙 \\
4 小时 & 是不是不太方便 \\
1 天 & 是不是不想理我 \\
\bottomrule
\end{manstable}

所以文字对话里，需要\textbf{显式地发回执}：

\begin{mdquote}
“我在看，你说。”

“刚忙完，这就回你。”

“你等我一下，我理一理再回。”
\end{mdquote}

这些话看起来是废话，信息量为零。\textbf{但它们的价值不在信息，在于把上面那张表里的时钟关掉。}它们的作用和“嗯”完全一样：告诉对方，链路还开着。

值得注意的是，文字消息有一个面对面和电话都没有的故障：\textbf{回执和内容会错位。}

\sample{0601}\par
\begin{mdquote}
周五晚上你在地铁上，给同事连着发了四条语音，每条四十秒。讲的是周一早会要用的那份材料怎么改：第一段是背景，第二段是具体改哪几页，第三段是格式要求，第四段是截止时间。

你发完就进了电梯。

六分钟后你出来，他的回复排了三条：

“第一条收到。”

“第二条在听。”

“等等，第三条怎么断在半句？”

你翻了一下才明白：\textbf{他在你发第二条的时候进了一个地下车库，信号断了。}第三条他压根没收到，只在电梯里缓存了一半。第四条要等他重新有信号才补上。

而他那边看到的顺序，和你这边的顺序，已经不是一回事了。
\end{mdquote}

这个样本要说明的是：\textbf{语音消息的“已发出”四个字，是发给你的，不是发给他的。}

你这边显示一个勾，只代表东西从你的设备出去了。它在路上发生了什么，你这边完全看不见。而人类会默认“已发出”等于“已收到”——\textbf{这条默认是错的，而且是在一个你根本不会怀疑它错的地方错的。}

所以用语音消息做确认，有一个额外的成本：\textbf{你发出的回执，本身也可能丢。}如果要紧的事，请回到同步信道（电话），或者要求对方对字面内容做一次显式复述。

\section{延伸：当对方是长辈或上级}

\textbf{当对方是长辈}：回执要稍微“重”一点。“嗯”在长辈听来可能显得敷衍，改用“我在听”“你说”这类完整的短句，成本几乎一样，效果好得多。

\textbf{当对方是上级}：回执除了告诉对方“我在听”，还要顺带证明“我在记”。所以可以加一句复述的引子：

\begin{mdquote}
“嗯，我记一下……您是说下周先不用报？”
\end{mdquote}

这一句同时完成了回执和确认理解，在职场里是更划算的用法。

\textbf{当对方是陌生人}：纯回执就够了。“嗯”在陌生人之间不带任何额外含义，是最安全的。

这里还有一类场景值得单列：\textbf{当对方看不见你的脸。}

\sample{0600}\par
\begin{mdquote}
周一下午三点，一条十六个人的视频会议。你在的那一段是别人在汇报数据，跟你没什么关系，你把摄像头关了，抓了张纸记了两笔。

会议开了四十分钟，你全程一个声音都没出。

四点零五分，主持人突然说了一句：“\textbf{刚才我没听错吧，好像有几位掉线了？}麻烦在线的同学扣个 1。”

聊天框里瞬间刷出十四行“1”。你是第十五个敲键盘的。

那一刻你才意识到：\textbf{在这四十分钟里，你和真正掉线的人，在屏幕上看起来一模一样。}
\end{mdquote}

视频会议是目前人类接口里最像机器的一种信道，也正因为像，它会\textbf{放大沉默的歧义}。

面对面时，你不出声，对方能看见你在点头。电话里，你不出声，对方至少知道你在那端（除非他怀疑信号）。视频会议里，\textbf{摄像头关着你是一片灰，摄像头开着但你在听报告，你也是一片灰。}信道把“沉默”和“离线”压缩成了同一个像素。

所以视频会议有一条额外的规则：

\textbf{在不是你发言的时段里，也要给出周期性的显式回执。}

具体说，就是在别人讲完一个段落的时候，对着麦克风说一句短的。一个“嗯”，一个“对”，或者干脆一句“我在”。成本接近零，但只要有一句，你就从“疑似掉线”的名单里出来了。

\section{文字消息的确认时延表}

这一节把 3.7 的那张时间曲线展开推演一遍，因为文字消息的误判，\textbf{大多数不是因为内容说错了，而是因为发送方在你沉默的每一分钟里，都在替你想一个更坏的理由。}

先给一个总的结论：\textbf{文字消息的传播延迟，不是网络延迟，是心理延迟。}

下面这条曲线，描述的是同一条消息——周三晚上九点发出去、对方还没回——在发送方心里走过的完整过程。

\begin{manstable}{P{2.6cm} P{2.6cm} L L}
\tbh{已经过去} & \tbh{发送方的心理状态} & \tbh{他通常对你做出的推断} & \tbh{他可能采取的行动} \\
\midrule
0—5 分钟 & 平静 & 在忙，正常 & 无 \\
5—30 分钟 & 略微留意 & 可能在开会／在开车 & 无，但记住了 \\
30 分钟—2 小时 & 开始回看自己的消息 & 是不是我说得不合适 & 翻一遍自己发的那段 \\
2—4 小时 & 转而下沉 & 是不是不太方便说 & 想撤回，但没有撤回按钮 \\
4—12 小时 & 怀疑对象从“事”转到“人” & 是不是不太想理我 & 开始给这件事找“他平时也这样”的证据 \\
12 小时—1 天 & 情绪转成低位 & 我是不是说错话了 & 下次不再主动找他 \\
超过 1 天 & 结案 & 他大概是不想理我 & 把你从“会主动联系的人”名单里移除 \\
\bottomrule
\end{manstable}

需要说明的是这张表里的几个关键点。

\textbf{第一，转折发生在 2 小时到 4 小时之间。}在那之前，发送方怀疑的是“客观条件”（你在忙）；在那之后，他怀疑的是“你的态度”。\textbf{这两件事的性质不一样，而且一旦跨过去，就很难拉回来。}

\textbf{第二，超过一天不回，不等于“对方生气”。}它更常见的后果是\textbf{对方以后不再主动}。这才是这条曲线的真正成本：\textbf{它消耗的不是这一次的耐心，是下一次的开口。}

\textbf{第三，也是最容易被忽略的：发送方会把你现在的沉默，和你以前的沉默加在一起算。}一个人如果经常几小时不回，偶尔一次秒回救不了他；反过来，一个平时回复很快的人，偶尔半天不回，对方会立刻察觉异常。\textbf{人类是拿历史记录来做当前判断的。}

\sample{0603}\par
\begin{mdquote}
周二中午十二点半，你给一个老朋友发了条消息：“下个月我可能去你那边出差，有空吃个饭吗。”

下午两点，他在开会。

下午五点，他在改方案，看到了没回，想着晚上再说。

晚上十点，他躺在床上刷手机，忽然想起来这条消息。他点开，打了一句“有啊”，删掉；改成“肯定有啊，你说哪天”，又觉得太用力；最后发出去的是——

“在的，你说。”

而你这边，从中午到晚上，已经把这件事在心里过了三遍。
\end{mdquote}

这个样本的重点不在他回复得晚，而在他最后发出去的话是\textbf{“在的，你说”}——\textbf{他其实是在替自己这九个小时的沉默发回执。}他在意的不是你等着急，而是他自己知道，这九个小时里你心里走过了一张什么样的表。

\section[常见误区]{常见误区}

\hcimistake{“我在听啊”当成回执}{\hciTag{反例}你讲了八分钟，对方突然说：“我在听啊。”}{它出现得太晚了。等到你不得不承认自己在听的时候，对方已经怀疑很久。回执要在\textbf{过程中}给，不在\textbf{结算时}给。}

\hcimistake{用眼神代替语言回执}{\hciTag{反例}会议桌上对方问：“这个方案你觉得怎么样？”你点了点头。}{面对面时通用，视频会议里勉强通用，电话和文字里等于零。点头的动作只在对方\textbf{能看见}你的时候才是回执。}

\hcimistake{“嗯”说成了“嗯——嗯”（拖长音）}{\hciTag{反例}“嗯——嗯——”}{在多数方言区，这个音是“敷衍”的意思，效果与回执相反。\textbf{回执的音长要短，拖长就变了性质。}}

\hcimistake{在对方讲到情绪最重的地方时保持沉默}{\hciTag{反例}对方说：“我妈住院那两天，我一个人在走廊坐了一整夜。”你什么都没说。}{这个位置最需要回执。沉默会让对方以为你在评判他。\textbf{情绪最重的句子后面，必须有一声“嗯”接住。}}

\hcimistake{把“行吧”当“好的”}{\hciTag{反例}“行吧。”——你理解为“同意”。}{\hciSee{3.3}。“行吧”是撤退，它的成本会延期结算。}

\hcimistake{群聊里用回执}{\hciTag{反例}同事在群里汇报进度，你回了一个“嗯”。}{一对一场景的规则不适用于群聊，群聊中每个人的回执会被理解为整群的态度。\textbf{你一个人的“嗯”，在群里看起来像是这个团队表示的。}}

\hcimistake{在对方已经明显不想说的时候强行要回执}{\hciTag{反例}“你怎么不说话？你到底怎么想的？”}{这句会让对方为沉默本身道歉，而不是继续讲。你要的是信息，但你逼出来的是\textbf{关于沉默的道歉}。}

\hcimistake{拿高强度的回执当礼貌}{\hciTag{反例}“对对对，你说得对，是是是，你说得都对。”}{\hciSee{3.6}。高强度的回执会被自动升级成同意，而且它同时拖低了所有回执的可信度——\textbf{给得太顺的回执，等于没有回执。}}

%% ---------- 本章小结 ----------
\hciblocktitle{bullseye}{本章小结}

综上所述，本章的核心结论可以概括为以下几点：

\begin{enumerate}[label=\bfseries\arabic*., leftmargin=2.4em, labelsep=0.7em,
                  itemsep=0.45em, topsep=0.2em, parsep=0.2em]
\item 沉默不等于在听。发送方看不到你的状态，只能靠声音判断。
\item “嗯”的作用不是回答，是回执：告诉对方链路还开着。
\item “嗯……”是预警，“行吧”是撤退。这两种最常被误读。
\item 确认的节奏大约在 20 到 40 秒一次，插在对方的自然停顿处；换气停顿不是窗口。
\item 确认过度的代价高于沉默：高强度回执会被自动升级成同意，也会让所有回执一起贬值。
\item 文字消息没有隐式回执，所以需要显式地发一句“我在看”——而且回执本身也可能在路上丢。
\item 文字消息的沉默有一个心理时延曲线：超过 2 到 4 小时，发送方的怀疑会从“你在忙”转向“你的态度”；超过一天，他不再主动。
\end{enumerate}

%% ---------- 练习 ----------
\begin{exercisessec}
\item 在接下来的一次通话中，刻意把回执频率控制在每 30 秒一次。对话结束后记录：对方有没有在某处放慢语速或停下来试探（那是他怀疑断线的信号）。

\item 找出你最近收到的三句“行吧”，回忆当时的上下文，判断它们是“同意”还是“撤退”。（[请在此处填写具体的三句话]）

\item 在你下一次要晚回一条消息的时候，先发一句“我在看，待会儿回你”。记录对方后续的回复与平时有没有不同。

\item 下一次和真人对谈时，把你自己的回执做一次降级：全程只用“嗯”，不用“对对对”“是是是”。结束后，请对方回答一个问题——\textbf{“刚才那段话里，你有哪一部分是觉得我已经同意了的？”}把他的回答记下来，看看你有没有在无意中承诺过什么。
\end{exercisessec}

%% ---------- 自测题 ----------
\begin{quizsec}
\item 下面哪一句是“撤退”，不是“同意”？

\begin{enumerate}[label=\Alph*., leftmargin=2.2em, labelsep=0.6em,
                  itemsep=0.2em, topsep=0.3em, parsep=0.1em]
\item 好的，那就这么办
\item 对，我同意
\item 行吧
\item 嗯嗯，听到了
\end{enumerate}

\item 判断：确认给得越多越安全，所以“嗯嗯嗯”比“嗯”更好。

\item 判断：对方沉默的时候，说明他在认真听。

\item 简答：视频会议里，为什么“全程不出声”和“真的掉线”会被混为一谈？

\item 简答：为什么文字消息需要显式地发一句“我在看”？
\end{quizsec}

%% ---------- 参考答案 ----------
\begin{answerssec}
\item C。“行吧”通常意味着“我说不过你，或者我不想吵了”。把它当成“好的”，会在很久以后收到一次你完全不理解的翻脸。\hciSee{3.3}。

\item 错。回执的频率该高，强度该低。“嗯嗯”这类高强度回执会被自动升级成同意，而且给得太顺的回执会一起贬值，反而更快被读成敷衍。\hciSee{3.6}。

\item 错。沉默和在听是两件长得一样的事。发送方看不到接收方的状态，只能靠声音判断。\hciSee{3.1}。

\item 因为视频会议这个信道，把“沉默”和“离线”压缩成了同一个像素：摄像头关着是一片灰，开着但在听也是一片灰。信道本身不携带“我在听”这个信息，所以必须靠显式的短句回执把两者分开。\hciSee{3.8}。

\item 因为文字通信没有语气、停顿、呼吸声这些隐式回执，发送方收不到任何返回，只能自己推测——而推测的方向通常越来越坏。这类句子信息量为零，作用是把那条“他是不是不想理我”的时钟关掉。\hciSee{3.7} 与 3.9。
\end{answerssec}

\hcirule

\begin{qabox}
\qaQ{读者提问}{如果对方一直讲，我插不进“嗯”，怎么办？}
\qaA{本手册回答}{这是一个非常好的问题。需要说明的是，插不进回执，通常说明对方也处在“只顾发送”的状态里。这种情况下可以放弃回执，改在对方换气的长停顿时补一句短的复述——它同时完成回执和理解确认，也更容易插进去。}
\end{qabox}

\hcirule

希望本章的内容能对你有所帮助。如需进一步了解第 4 章《信噪比》的相关内容，我可以随时为你展开。

%% 章末「页脚交互条」由 hci-manual.cls 自动挂接，无需手写。
